% Revisión académica e industrial, formato mínimo: una figura, una tabla y un
% párrafo por revisión. El resto de lo construido (tabla de protocolo,
% capas del corpus, auditoría red-team, normas, tendencias patentarias,
% secuencia funcional) se traslada íntegro al Anexo~\ref{app:review-full-detail}.
\subsection{Revisión académica e industrial}
\label{subsec:tf-discriminant}

La revisión se organiza en dos frentes con criterios de inclusión distintos. El
académico delimita qué mecanismos de decisión están publicados y con qué grado
de explicitud; el industrial delimita qué de eso llega a producto o a norma. Los
dos protocolos completos, con consultas, fechas de corte y criterios de
exclusión, están en el Anexo~\ref{app:review-full-detail}.

\paragraph{Académica}
Se revisaron 59 trabajos canónicos de coalición/MRTA y transporte cooperativo,
recuperados de Crossref, OpenAlex y \emph{snowballing} (protocolo completo en
el Anexo~\ref{app:review-full-detail}); 32 de ellos, con coordenada
discriminante completa en la fuente, se agrupan en la
Figura~\ref{fig:lit-methodological-map} por autoridad de decisión y
explicitud del mecanismo. La Figura~\ref{fig:lit-capability-matrix} cruza los
doce trabajos más próximos contra ocho capacidades y evidencia que ninguno
reúne certificación física, ejecución distribuida y reemplazo a la vez. La
auditoría adversarial de novedad contra los seis antecedentes más directos
del corpus completo se detalla en el Anexo~\ref{app:review-full-detail}.

% Figura 1 del consolidado académico: mapa metodológico de los 32 trabajos
% discriminantes del núcleo canónico. El cuerpo TikZ y los macros \mappointL y
% \mappointR se portan desde `thesis-reviews/academic_literature_review (1).tex`
% (líneas 60-148), cuyo conjunto de trabajos coincide con el del PDF vigente.
\begin{figure}[!htbp]
  \centering
  \makebox[\textwidth][c]{\resizebox{0.98\textwidth}{!}{%
\begin{tikzpicture}[x=80mm,y=54mm,font=\sffamily]
  \fill[softblue] (-1,-1) rectangle (0,0);
  \fill[softgreen] (0,-1) rectangle (1,0);
  \fill[softplum!55!white] (-1,0) rectangle (0,1);
  \fill[softwarm] (0,0) rectangle (1,1);
  \foreach \x in {-1,-.75,-.5,-.25,0,.25,.5,.75,1} \draw[rule,line width=.28pt] (\x,-1)--(\x,1);
  \foreach \y in {-1,-.75,-.5,-.25,0,.25,.5,.75,1} \draw[rule,line width=.28pt] (-1,\y)--(1,\y);
  \draw[-{Stealth[length=2.3mm]},ink,line width=.68pt] (-1.04,0)--(1.075,0);
  \draw[-{Stealth[length=2.3mm]},ink,line width=.68pt] (0,-1.04)--(0,1.09);
  \node[anchor=south,font=\bfseries\fontsize{7.2}{7.5}\selectfont] at (0,1.10) {data-driven};
  \node[anchor=north,font=\bfseries\fontsize{7.2}{7.5}\selectfont] at (0,-1.09) {white-box};
  \node[rotate=90,anchor=south,font=\bfseries\fontsize{7}{7.3}\selectfont] at (-1.12,0) {descentralizado};
  \node[rotate=90,anchor=north,font=\bfseries\fontsize{7}{7.3}\selectfont] at (1.12,0) {centralizado};
  \node[align=center,font=\bfseries\fontsize{6.35}{6.65}\selectfont] at (-.60,.90) {distribuido\\[-.35mm]data-driven};
  \node[align=center,font=\bfseries\fontsize{6.35}{6.65}\selectfont] at (.60,.90) {centralizado\\[-.35mm]data-driven};
  \node[align=center,font=\bfseries\fontsize{6.35}{6.65}\selectfont] at (-.60,-.90) {distribuido\\[-.35mm]white-box};
  \node[align=center,font=\bfseries\fontsize{6.35}{6.65}\selectfont] at (.60,-.90) {centralizado\\[-.35mm]white-box};

  % data-driven, distributed / mixed
  \mappointL{-0.82}{0.72}{Bezerra}{2025}{B}{20}{learncol}
  \mappointL{-0.64}{0.55}{Huo}{2025}{H}{19}{learncol}
  \mappointL{-0.46}{0.73}{Dai}{2024}{D}{17}{learncol}
  \mappointL{-0.78}{0.35}{Shibata-RL}{2023}{S}{16}{learncol}
  \mappointL{-0.43}{0.32}{Pandit}{2026}{P}{28}{learncol}
  \mappointR{-0.13}{0.50}{Eoh/Park}{2021}{E}{29}{learncol}

  % data-driven, centralized/mixed
  \mappointL{0.34}{0.64}{SADCHER}{2026}{S}{25}{learncol}
  \mappointL{0.62}{0.42}{CF-HMRTA}{2025}{C}{21}{searchcol}

  % white-box distributed - staggered
  \mappointL{-0.90}{-0.18}{Pereira}{2004}{P}{1}{controlcol}
  \mappointL{-0.78}{-0.34}{Parker/Tang}{2006}{PT}{2}{foundcol}
  \mappointL{-0.88}{-0.54}{Vig/Adams}{2006}{V}{3}{gamecol}
  \mappointL{-0.69}{-0.17}{CBBA}{2009}{C}{4}{gamecol}
  \mappointL{-0.64}{-0.39}{Service}{2011}{S}{5}{foundcol}
  \mappointL{-0.57}{-0.60}{Chen/Sun}{2012}{CS}{6}{gamecol}
  \mappointL{-0.45}{-0.22}{Jang}{2018}{J}{8}{gamecol}
  \mappointL{-0.39}{-0.47}{Dutta}{2019}{D}{10}{gamecol}
  \mappointL{-0.36}{-0.67}{Dutta}{2021}{D}{12}{gamecol}
  \mappointL{-0.14}{-0.28}{Ebel}{2024}{E}{18}{controlcol}
  \mappointR{-0.03}{-0.50}{Shibata}{2023}{S}{14}{controlcol}
  \mappointR{-0.04}{-0.71}{Fukao}{2025}{F}{22}{controlcol}
  \mappointL{-0.48}{-0.90}{GRAPE-S}{2026}{G}{24}{gamecol}

  % white-box central/mixed
  \mappointL{0.08}{-0.18}{Alonso-Mora}{2017}{AM}{7}{controlcol}
  \mappointL{0.20}{-0.40}{Machado}{2019}{M}{9}{controlcol}
  \mappointL{0.31}{-0.65}{Verma repl.}{2019}{V}{11}{controlcol}
  \mappointL{0.47}{-0.22}{Elaamery}{2021}{E}{13}{controlcol}
  \mappointL{0.59}{-0.44}{Gong}{2023}{G}{15}{controlcol}
  \mappointL{0.76}{-0.22}{Rizzo}{2020}{R}{30}{controlcol}
  \mappointR{0.89}{-0.40}{Loh}{2012}{L}{31}{controlcol}
  \mappointR{0.83}{-0.64}{Yufka}{2015}{Y}{32}{controlcol}
  \mappointR{0.61}{-0.73}{Muhammed}{2026}{M}{26}{controlcol}
  \mappointR{0.34}{-0.88}{Sahu}{2026}{S}{27}{controlcol}
  \mappointL{0.02}{-0.90}{Tian}{2025}{T}{23}{controlcol}

  \node[star,star points=5,star point ratio=2.1,fill=orange,draw=white,line width=.8pt,minimum size=8.5mm] at (-0.18,-0.78) {};
  \node[anchor=west,font=\bfseries\fontsize{7.0}{7.4}\selectfont,align=left] at (-0.09,-0.75) {Este TFM\\integración};
\end{tikzpicture}%
  }}
  \par\vspace{1mm}
  {\centering\fontsize{8.35}{8.9}\selectfont
  \chip{gamecol}{juegos/mercados}\hspace{5.5mm}
  \chip{controlcol}{control/consenso}\hspace{5.5mm}
  \chip{searchcol}{búsqueda/heurística}\hspace{5.5mm}
  \chip{learncol}{política aprendida}\hspace{5.5mm}
  \chip{foundcol}{fundacional}\par}
  \caption[Mapa metodológico del núcleo canónico.]{Mapa metodológico de 32
  trabajos discriminantes del núcleo canónico. El eje horizontal codifica
  autoridad de decisión (descentralizado $\rightarrow$ centralizado) y el
  vertical la explicitud del mecanismo (white-box $\rightarrow$ data-driven); el
  color identifica la familia primaria. Los trabajos de coalición y MRTA se
  concentran en la decisión lógica y los de transporte cooperativo en la
  ejecución física. Las posiciones son ordinales y se desplazan solo para evitar
  solapes: no representan calidad ni rendimiento.}
  \label{fig:lit-methodological-map}
  \viusource{elaboración propia a partir de la codificación del núcleo canónico}
\end{figure}

% Figura 2(a) y 2(b) del consolidado académico. La matriz se porta desde
% `thesis-reviews/academic_literature_review (1).tex` (líneas 176-217) y se le
% añade la columna C ---nivel de centralización C0--C4--- que incorpora el PDF
% vigente. Las modalidades físicas conservan los rótulos del original.
\begin{figure}[!htbp]
  \centering
  \makebox[\textwidth][c]{\resizebox{0.99\textwidth}{!}{%
\begin{minipage}{\linewidth}
\renewcommand{\arraystretch}{1.45}\setlength{\tabcolsep}{2.25pt}\fontsize{7.55}{8.0}\selectfont
\begin{tabularx}{\linewidth}{J{27mm}*{8}{K{10.75mm}}K{9.6mm}K{9.2mm}}
\toprule
\textbf{Trabajo} & \rotatebox{90}{\textbf{heterog.}} & \rotatebox{90}{\textbf{coal. var.}} & \rotatebox{90}{\textbf{reclut. local}} & \rotatebox{90}{\textbf{shared-load}} & \rotatebox{90}{\textbf{cert. física}} & \rotatebox{90}{\textbf{exec. dist.}} & \rotatebox{90}{\textbf{recovery}} & \rotatebox{90}{\textbf{replacement}} & \textbf{Ev.} & \textbf{C}\\
\midrule
Jang 2018 [8] & \capno & \capyes & \capyes & \capno & \capno & \capyes & \capno & \capno & T+S & C4\\
Dutta 2021 [12] & \capyes & \capyes & \capyes & \capno & \capno & \capyes & \capno & \capno & T+S & C4\\
GRAPE-S 2026 [24] & \capyes & \capyes & \cappart & \capno & \capno & \capyes & \capno & \capno & T+S & C3--4\\
CF-HMRTA 2025 [21] & \capyes & \capyes & \capno & \capno & \capno & \capno & \capno & \capno & T+S & C0--1\\
Huo 2025 [19] & \capyes & \capyes & \capyes & \capno & \capno & \capyes & \cappart & \capno & T+S & C3--4\\
Bezerra 2025 [20] & \capyes & \capyes & \capyes & \capno & \capno & \capyes & \cappart & \capno & S & C4\\
Alonso-Mora 2017 [7] & \capno & \cappart & \capno & \capyes & \capyes & \cappart & \cappart & \capno & T+S+P & C0--1\\
Ebel 2024 [18] & \capno & \capno & \capno & \capyes & \capyes & \capyes & \capno & \capno & T+S+P & C3\\
Shibata 2023 [14] & \capno & \cappart & \capno & \capyes & \capyes & \capyes & \cappart & \cappart & S+P & C4\\
Muhammed 2026 [26] & \capno & \capno & \capno & \capyes & \capyes & \capyes & \capno & \capno & T+S+P & C3\\
Verma repl. 2019 [11] & \capno & \cappart & \capno & \capyes & \cappart & \capno & \capyes & \capyes & S+P & C0--1\\
Sahu 2026 [27] & \capno & \cappart & \capno & \capyes & \cappart & \cappart & \capyes & \capyes & S+P & C1--2\\
\bottomrule
\end{tabularx}
\end{minipage}%
  }}
  \par\vspace{2mm}
  \makebox[\textwidth][c]{\resizebox{0.92\textwidth}{!}{%
\begin{tikzpicture}[x=1mm,y=1mm,font=\sffamily\fontsize{6.25}{6.65}\selectfont]
\tikzset{bx/.style={rounded corners=1pt,draw=rule,line width=.35pt,minimum width=41.5mm,minimum height=17mm,align=center}}
\node[bx,fill=softblue](s)at(21,10){\textbf{Supported cargo}\\carga soportada\\wrench y límites};
\node[bx,fill=softwarm](r)at(65,10){\textbf{Rigid attachment}\\formación rígida\\distribución de fuerza};
\node[bx,fill=softgreen](p)at(109,10){\textbf{Pushing / caging}\\contacto unilateral\\geometría no-prehensil};
\node[bx,fill=softplum](g)at(153,10){\textbf{Grasping / arms}\\grasp matrix\\force closure};
\foreach \a/\b in {s/r,r/p,p/g}{\draw[-{Stealth[length=1.55mm]},muted,line width=.42pt](\a)--(\b);}
\end{tikzpicture}%
  }}
  \caption[Matriz de capacidades y modalidades físicas.]{(a) Matriz de
  capacidades de los trabajos discriminantes: \capyes\ evidencia explícita,
  \cappart\ cobertura parcial o condicionada y \capno\ no documentado en el
  alcance revisado. La columna Ev. indica el tipo de evidencia (T teórica,
  S simulación, P plataforma física) y la columna C el nivel de centralización,
  de C0 ---decisión centralizada--- a C4 ---decisión local con intercambios
  vecinales sin autoridad central en tiempo de ejecución---. (b) Modalidades
  físicas que no deben mezclarse bajo la etiqueta genérica \emph{cooperative
  transport}: cada una impone un conjunto de esfuerzos admisibles distinto.}
  \label{fig:lit-capability-matrix}
  \viusource{elaboración propia a partir de la codificación del núcleo canónico}
\end{figure}


\begin{tcolorbox}[enhanced,breakable,colback=VIUOrange!5,colframe=VIUOrange,
  boxrule=0.45pt,leftrule=2.8pt,arc=1.2mm,left=2.5mm,right=2.5mm,top=1.8mm,bottom=1.8mm]
\textbf{Brecha académica defendible.} En el corpus revisado no se identificó una
arquitectura que integre, en una misma cadena operativa, la selección local de una
coalición física heterogénea de tamaño variable desde una flota mayor, la verificación
mecánica de factibilidad de la carga compartida y la recomposición en línea de un
miembro durante el transporte, manteniendo la decisión de coalición y la ejecución sin
una autoridad global permanente --- la misma combinación que la lectura de la matriz
industrial (Tabla~\ref{tab:tf-industrial-matrix}) tampoco encuentra resuelta a la vez. El
enunciado se restringe al protocolo de búsqueda documentado y no afirma inexistencia
absoluta.
\end{tcolorbox}

\paragraph{Industrial}
Se auditaron catorce precedentes comerciales con fuente oficial (protocolo en el
Anexo~\ref{app:review-full-detail}), situados en la
Figura~\ref{fig:tf-industrial-map} por interacción cooperativa y
descentralización de la coordinación. La Tabla~\ref{tab:tf-industrial-matrix}
muestra que la industria resuelve por separado carga compartida, emparejamiento
con relevo ante fallo y coordinación de flota, pero ningún precedente integra las
tres a la vez con selección local y certificación mecánica. El marco
normativo aplicable (VDA~5050, ISO~21423, ISO~3691-4, ANSI/A3~R15.08-3) no
cubre esa combinación, según se detalla en el Anexo~\ref{app:review-full-detail}.

% Figura 1 del consolidado industrial, portada literalmente desde
% `thesis-reviews/TFM_estado_industrial_patentes_EDITABLE_v47.tex` (líneas
% 148-247). Solo se cambia el envoltorio: el original usa \figcap y aquí va
% dentro de un `figure` con el pie y la fuente en el estilo de la memoria.
% El cuerpo TikZ y la paleta son los del autor, de modo que sigue siendo
% editable con la misma sintaxis que en el consolidado.
\begin{figure}[!htbp]
  \centering
  \makebox[\textwidth][c]{\resizebox{\textwidth}{!}{%
\begin{tikzpicture}[x=1mm,y=1mm,font=\sffamily\fontsize{6.45}{6.95}\selectfont]
% La caja original (0,-14)--(174,98) dejaba fuera las etiquetas ancladas al este
% del borde derecho, que al reescalar invadían el margen. Se amplía para que
% \resizebox mida el contenido completo.
\path[use as bounding box] (0,-14) rectangle (196,98);

% plot frame and regions: four clear quadrants
\fill[gray!8] (30,16) rectangle (98,53);
\fill[orange!10] (98,16) rectangle (166,53);
\fill[teal!10] (30,53) rectangle (98,90);
\fill[navy!8] (98,53) rectangle (166,90);
\draw[rule,line width=.45pt] (30,16) rectangle (166,90);
\draw[rule!88,line width=.42pt] (98,16)--(98,90);
\draw[rule!88,line width=.42pt] (30,53)--(166,53);
\foreach \x in {64,132}{\draw[rule!75,densely dotted,line width=.24pt] (\x,16)--(\x,90);}
\foreach \y in {34.5,71.5}{\draw[rule!75,densely dotted,line width=.24pt] (30,\y)--(166,\y);}

% axes
\node[anchor=north,align=center,text width=30mm] at (47,14.7){Flota\ independiente};
\node[anchor=north,align=center,text width=30mm] at (81,14.7){Movimiento\\ coordinado};
\node[anchor=north,align=center,text width=30mm] at (115,14.7){Misma carga\\ sincronizada};
\node[anchor=north,align=center,text width=30mm] at (149,14.7){Acoplamiento\\ rígido/modular};
\node[anchor=north,font=\bfseries] at (98,5.2){Interacción cooperativa entre vehículos};
\node[anchor=north,font=\sffamily\fontsize{5.5}{5.8}\selectfont,text=muted] at (98,9.4){menor $\rightarrow$ mayor};

\node[anchor=east,align=right,text width=26mm] at (27,25){Humano /\\ configurado};
\node[anchor=east,align=right,text width=26mm] at (27,43.5){Control\\ central};
\node[anchor=east,align=right,text width=26mm] at (27,62){Líder /\\ master};
\node[anchor=east,align=right,text width=26mm] at (27,80.5){Entre vehículos /\\ distribuido};
\node[rotate=90,anchor=south,font=\bfseries] at (4.5,53){Descentralización de la coordinación y el control};
\node[rotate=90,anchor=south,font=\sffamily\fontsize{5.3}{5.6}\selectfont,text=muted] at (9.6,53){menor $\rightarrow$ mayor};

% label and point styles
\tikzset{regtitle/.style={anchor=west,font=\bfseries,fill=white,fill opacity=.84,text opacity=1,inner sep=1.1pt,rounded corners=.8pt}, regsub/.style={anchor=west,text=muted,fill=white,fill opacity=.84,text opacity=1,inner sep=.9pt,rounded corners=.8pt}, pt/.style={circle,draw=onyx,line width=.55pt,minimum size=5.2mm,inner sep=0pt,fill=white}, num/.style={font=\sffamily\fontsize{5.8}{6.0}\selectfont\bfseries,text=white}, ctx/.style={diamond,aspect=1.15,draw=onyx!65,line width=.45pt,minimum width=5.0mm,minimum height=4.3mm,inner sep=0pt,fill=onyx!14!white}, ctxnum/.style={font=\sffamily\fontsize{5.1}{5.3}\selectfont\bfseries,text=onyx}}

% region labels
\node[regtitle,text=teal!90!black] at (34,87.8){flota autónoma};
\node[regsub] at (34,84.4){coordinación distribuida; sin carga compartida};
\node[regtitle,text=navy] at (101,87.8){convergencia objetivo};
\node[regsub] at (101,84.4){carga compartida + mayor autonomía};
\node[regtitle,text=ink] at (34,20.8){automatización convencional};
\node[regsub] at (34,17.6){1 carga por robot; control configurado o central};
\node[regtitle,text=orange!90!black] at (100.5,18.9){heavy-load configurado};
\node[regsub] at (100.5,15.9){carga compartida/acople; baja autonomía};

% points + labels
\node[pt,fill=orange] (c1) at (151,25) {}; \node[num] at (c1){1};
\draw[orange!70,line width=.32pt] (c1)--(141.5,23.3); \node[anchor=east,align=right] at (141.0,23.3){KUKA omniMove / Airbus\\2016};

\node[pt,fill=teal] (c2) at (136.5,79.0) {}; \node[num] at (c2){2};
\draw[teal!70,line width=.32pt] (c2)--(128.8,75.4); \node[anchor=east,align=right] at (128.3,75.4){Stäubli + R3\\2022};

\node[pt,fill=blue] (c3) at (116,47) {}; \node[num,text=onyx] at (c3){3};
\draw[blue!70,line width=.32pt] (c3)--(109.5,55.5); \node[anchor=east,align=right] at (109.0,55.5){SIASUN dual-vehicle\\2024};

\node[pt,fill=spicy!68!white] (c4) at (113,63) {}; \node[num] at (c4){4};
\draw[plum!70,line width=.32pt] (c4)--(103,67); \node[anchor=east,align=right] at (102.5,67){KUKA KMP 3000P\\2025};

\node[pt,fill=orange] (c5) at (136,29.5) {}; \node[num] at (c5){5};
\draw[orange!70,line width=.32pt] (c5)--(127.5,28.8); \node[anchor=east,align=right] at (127.0,28.8){GREATOX GTVL25\\cat. 2026};

\node[pt,fill=orange] (c6) at (158,32) {}; \node[num] at (c6){6};
\draw[orange!70,line width=.32pt] (c6)--(165,39.0); \node[anchor=west,align=left] at (165.5,39.0){TII SCHEUERLE SPMT\\2023};

\node[pt,fill=teal] (c7) at (48,83) {}; \node[num] at (c7){7};
\draw[teal!70,line width=.32pt] (c7)--(56,78.5); \node[anchor=west,align=left] at (56.5,78.5){AGILOX X-SWARM / BMW\\2023};

\node[pt,fill=spicy!68!white] (c8) at (131,64) {}; \node[num] at (c8){8};
\draw[spicy!70,line width=.32pt] (c8)--(140.0,59.0); \node[anchor=west,align=left] at (140.5,59.0){Beacon Dual-AMR\\2026};

\node[pt,fill=orange] (c9) at (148,38) {}; \node[num] at (c9){9};
\draw[orange!70,line width=.32pt] (c9)--(137.5,47.0); \node[anchor=east,align=right] at (137.0,47.0){HUBTEX Aviation\\2020};

\node[pt,fill=blue] (c10) at (53,51.5) {}; \node[num,text=onyx] at (c10){10};
\draw[blue!70,line width=.32pt] (c10)--(60,54); \node[anchor=west,align=left] at (60.5,54){Hyundai HMGICS\\2023};

% contextual central-fleet systems (not expanded in Table 1)
\node[ctx] (ctxA) at (40.5,42.0) {}; \node[ctxnum] at (ctxA){A};
\draw[mustard!80!onyx,line width=.28pt] (ctxA)--(35.0,38.5); \node[anchor=east,align=right,font=\sffamily\fontsize{5.7}{6.0}\selectfont] at (34.5,38.5){MiR Fleet};

\node[ctx] (ctxB) at (58.0,39.0) {}; \node[ctxnum] at (ctxB){B};
\draw[mustard!80!onyx,line width=.28pt] (ctxB)--(61.5,35.0); \node[anchor=west,align=left,font=\sffamily\fontsize{5.7}{6.0}\selectfont] at (62.0,35.0){OTTO Fleet\\Manager};

\node[ctx] (ctxC) at (74.0,46.0) {}; \node[ctxnum] at (ctxC){C};
\draw[mustard!80!onyx,line width=.28pt] (ctxC)--(77.0,41.8); \node[anchor=north,align=center,font=\sffamily\fontsize{5.7}{6.0}\selectfont] at (77.0,41.4){Geek+ RMS};

\node[ctx] (ctxD) at (89.0,40.5) {}; \node[ctxnum] at (ctxD){D};
\draw[mustard!80!onyx,line width=.28pt] (ctxD)--(92.0,36.5); \node[anchor=west,align=left,font=\sffamily\fontsize{5.7}{6.0}\selectfont] at (92.5,36.5){Swisslog\\IntraMove};

% convergence vectors and TFM target
% The target is not placed at the rigid-coupling extreme: the TFM seeks carga compartida
% multi-contact cooperation with distributed coalition decisions, not necessarily a rigid robot-robot module.
\node[rounded corners=1.0pt,minimum size=6.4mm,inner sep=0pt,draw=spicy,line width=.55pt,fill=autumn] (tfm) at (147.5,84.0) {};
\node[anchor=west,font=\bfseries,text=spicy] at (152.2,84.0){TFM objetivo};

% compact legend
\node[anchor=north west,align=left,text width=162mm,text=muted] at (5,1.5){\fontsize{6.4}{7}\selectfont Forma: círculo = caso principal; rombo gris = contexto de flota; cuadrado naranja = TFM.\\ \textbf{En los círculos}, el color indica: \textcolor{orange}{\textbf{configurado/humano}}, \textcolor{blue}{\textbf{central}}, \textcolor{plum}{\textbf{líder/seguidor}}, \textcolor{teal}{\textbf{entre vehículos/distribuido}}.};
\end{tikzpicture}%
  }}
  \caption[Precedentes industriales de cooperación multi-vehículo.]{Diez
  precedentes principales y cuatro referencias contextuales de gestión de flota,
  clasificados según dos dimensiones ordinales: interacción cooperativa entre
  vehículos (izquierda $\rightarrow$ derecha) y descentralización de la
  coordinación y el control (abajo $\rightarrow$ arriba). Los círculos numerados
  corresponden a la Tabla~\ref{tab:tf-industrial-matrix}; los rombos A--D son
  contexto de flota centralizada. La posición dentro de una categoría se desplaza
  solo para evitar solapes y no tiene interpretación métrica.}
  \label{fig:tf-industrial-map}
  \viusource{elaboración propia a partir de las fuentes oficiales citadas en la
  Tabla~\ref{tab:tf-industrial-matrix}}
\end{figure}

% Tabla 1 del consolidado industrial, portada desde
% `thesis-reviews/TFM_estado_industrial_patentes_EDITABLE_v47.tex` (líneas
% 263-334). Cambios respecto del original, todos de integración:
%   - los tipos de columna L{}/C{}/Y pasan a J{}/K{}/Z, porque C e Y ya
%     existen en viu-mrob-thesis.sty con otra signatura;
%   - el hueco de fotografía se sustituye por el número de caso, que es lo
%     que casa la fila con la Figura del mapa; el enlace a la fuente oficial
%     se conserva sobre el nombre del sistema;
%   - las marcas \casecite se sustituyen por citas biblatex reales.
\begin{table}[!htbp]
  \centering
  \caption{Matriz funcional de precedentes industriales frente a la cadena operativa del trabajo.}
  \label{tab:tf-industrial-matrix}
{\setstretch{1}\fontsize{7.35}{8.10}\selectfont
 \setlength{\tabcolsep}{0.4pt}
 \renewcommand{\arraystretch}{1.32}
\begin{tabularx}{\linewidth}{@{}J{3.72cm}K{0.74cm}K{0.78cm}K{0.86cm}K{0.86cm}K{0.82cm}K{0.90cm}Z@{}}
\toprule
\textbf{Sistema / evidencia} &
{\fontsize{4.9}{5.2}\selectfont\shortstack{\textbf{Carga}\\\textbf{comp.}}} &
{\fontsize{4.9}{5.2}\selectfont\shortstack{\textbf{Multi-veh.}\\\textbf{($>2$)}}} &
{\fontsize{4.9}{5.2}\selectfont\shortstack{\textbf{Selecc.}\\\textbf{miembros}}} &
{\fontsize{4.9}{5.2}\selectfont\shortstack{\textbf{Coord./ctrl.}\\\textbf{local}}} &
{\fontsize{4.9}{5.2}\selectfont\shortstack{\textbf{Cert.}\\\textbf{factib.}}} &
{\fontsize{4.9}{5.2}\selectfont\shortstack{\textbf{Recup.}\\\textbf{coalición}}} &
\textbf{Evidencia industrial / límite}\\
\midrule
\rowcolor{soft}
\matrixsystem{1}{KUKA · omniMove}{Airbus Stade}{2016}{https://www.kuka.com/en-us/industries/solutions-database/2016/07/solution-robotics-airbus}
& \capyes & \capunk & \capno & \capno & \capno & \capno
& 100 t; composición dedicada. Sin selección dinámica, certificado online ni recomposición autónoma publicada \parencite{kukaOmnimoveAirbus2016}.\\

\matrixsystem{2}{Stäubli + R3}{Virtual drawbar}{2022}{https://www.r3.group/en/success-stories/collaborative-agvs-in-the-cleanroom}
& \capyes & \capno & \capno & \cappart & \capno & \capno
& 2 AGV, 2$\times$5 t; ejecución V2V directa a 32 ms con PROFINET/PROFIsafe. La fuente no demuestra formación distribuida de la pareja ni ausencia de una autoridad dominante \parencite{staubliR3Drawbar2022}.\\

\rowcolor{soft}
\matrixsystem{3}{SIASUN}{Heavy-truck dual vehicle}{2024}{https://en.siasun.com/heavy-truck-assembly-line-mobile-robot.html}
& \capyes & \capno & \cappart & \capno & \capno & \capno
& Dual-vehicle linkage + MES. Selección a nivel de planta, no reclutamiento local ni recomposición física \parencite{siasunDualVehicle2024}.\\

\matrixsystem{4}{KUKA · KMP 3000P}{Load sharing}{2025}{https://www.kuka.com/en-us/products/amr-autonomous-mobile-robotics/mobile-platforms/kmp-3000p-omnimove}
& \capyes & \capno & \capno & \cappart & \capno & \capno
& Dos plataformas, hasta 6 t; líder--seguidor. Carga compartida comercial, pero pareja fija \parencite{kukaKmp3000p2025}.\\

\rowcolor{soft}
\matrixsystem{5}{GREATOX · GTVL25}{Multi-vehicle synchronous}{cat. 2026}{https://en.greatoxagv.com/products_info/id-14.html}
& \capyes & \cappart & \capno & \capunk & \capno & \capunk
& 80--220 t y cooperación multi-vehículo. Monitoriza centro de gravedad; no publica certificado online de factibilidad mecánica, selección o recuperación \parencite{greatoxGtvl2026}.\\

\matrixsystem{6}{TII SCHEUERLE · SPMT}{Modular heavy transport}{2023}{https://www.tii-group.com/tii-scheuerle/our-solutions/spmt/scheuerle-spmt}
& \capyes & \cappart & \capno & \capno & \capno & \capno
& \textit{Project cargo} modular; configuración de ingeniería, no certificado online ni reclutamiento autónomo \parencite{tiiScheuerleSpmt2023}.\\

\rowcolor{soft}
\matrixsystem{7}{AGILOX · X-SWARM}{BMW Regensburg}{2023}{https://www.agilox.net/en/blog/bmw-agilox-swarm/}
& \capno & \capna & \capna & \capyes & \capna & \capna
& 27 AMR; coordinación entre pares sin líder central. Una carga por AMR: no aplica coalición física \parencite{agiloxXswarmBmw2023}.\\

\matrixsystem{8}{Beacon Robot · Dual-AMR}{Collaborative transport}{2026}{https://www.beacon-robot.com/news/beacon-robot-amr-dual-robot-collaboration-officially-released-from-working-alone-to-working-together/}
& \capyes & \capno & \capyes & \cappart & \capno & \capno
& Emparejamiento automático y relevo de autoridad ante fallo documentados. 2 robots; despacho central y líder--seguidor. No se documenta sustitución/reclutamiento de un nuevo miembro \parencite{beaconDualAmr2026}.\\

\rowcolor{soft}
\matrixsystem{9}{HUBTEX · SL/HL-AGV}{Aviation}{2020}{https://www.hubtex.com/sites/default/files/2020-04/HUBTEX_Aircraft\%20Components_ENG.pdf}
& \capyes & \capunk & \capno & \capno & \capno & \capno
& Hasta 70 t; plataformas acoplables y operación semi/remota. Sin selección dinámica, certificado online o recomposición \parencite{hubtexAviation2020}.\\

\matrixsystem{10}{Hyundai · HMGICS}{AGV / AMR}{2023}{https://www.hyundaimotorgroup.com/en/story/CONT0000000000126181}
& \capno & \capna & \capna & \cappart & \capna & \capna
& Orquestación de planta y reactividad local; sin coalición física de carga compartida \parencite{hyundaiHmgics2023}.\\
\midrule
\rowcolor{softplum}
\multicolumn{1}{@{}l}{\parbox[t]{3.55cm}{\textbf{TFM objetivo}\\[-0.3pt]{\fontsize{5.35}{5.65}\selectfont\color{muted}(referencia conceptual)}}}
& \capgoal & \capgoal & \capgoal & \capgoal & \capgoal & \capgoal
& Integrar carga compartida, coalición heterogénea de tamaño adaptable, reclutamiento local, coordinación/control distribuido, verificación de factibilidad contacto/torsor y recomposición tras fallo.\\
\bottomrule
\end{tabularx}}

\vspace{2.6pt}
\noindent\colorbox{softwarm}{\parbox{0.972\linewidth}{\setstretch{1}\fontsize{7.35}{8.1}\selectfont
\textbf{Lectura de la matriz.} La industria demuestra por separado carga compartida, emparejamiento/relevo ante fallo y coordinación distribuida de flota. \textbf{Ninguno de los precedentes revisados documenta simultáneamente una coalición física heterogénea multi-vehículo, selección local de miembros, verificación mecánica de factibilidad y recomposición durante la misión.}}}

\vspace{1.8pt}
{\setstretch{1}\fontsize{7.15}{8.0}\selectfont\color{muted}
Contexto A--D del mapa: MiR Fleet, OTTO Fleet Manager, Geek+ RMS y Swisslog IntraMove coordinan flotas de carga individual; no se incluyen como filas principales porque no aportan carga compartida. Los símbolos indican evidencia pública revisada, no una puntuación de similitud. En Multi-veh. ($>2$), \cappart\ indica cooperación multi-vehículo documentada sin cardinalidad superior a dos explícita en la fuente.\par}

\vspace{2pt}
{\centering\footnotesize
\capyes\ documentado \quad \cappart\ parcial o condicionado \quad
\capno\ no documentado públicamente \quad \capunk\ arquitectura no publicada \quad
\capna\ no aplica \quad \capgoal\ objetivo del TFM.\\
«No documentado» no equivale a inexistencia.\par}
  \viusource{elaboración propia con documentación técnica oficial de cada fabricante}
\end{table}

